\documentclass{article}
\usepackage{import}
\import{../../template/}{article-template}
\graphicspath{{../../images/}}
\addbibresource{../../template/library.bib}

\title{Shell}
\author{PENG}
\date{\today}

\begin{document}
\maketitle
\tableofcontents
\section{Introduction}
\url{https://www.gnu.org/software/bash/manual/bash.html}

\section{Shell Syntax}
\subsection{Quoting}
\begin{definition}[quoting]
    Quoting is used to remove the special meaning of certain characters or words to the shell.
    \begin{verbatim}
echo "${var}"
echo "$(date)"
echo "`date`"
echo -e "\t"
# \n : newline
# \t : horizontal tab
# \r : carriage return        
\end{verbatim}
\end{definition}

\subsection{Pipelines}
\begin{definition}[pipeline]
    cmd1 output $\to$ cmd2 input
    \begin{verbatim}
cmd1 | cmd2      # stdout 
cmd1 |& cmd2     # stderr to stdout
cmd1 2>&1 | cmd2 # stderr to stdout        
\end{verbatim}
\end{definition}

\subsection{Lists of Commands}
\begin{definition}[list]~
    \begin{verbatim}
cmd1 && cmd2 # And, cmd1 true then cmd2
cmd1 || cmd2 # Or, cmd1 false then cmd2
cmd1; cmd2  
cmd1 & cmd2  # cmd1 async
\end{verbatim}
\end{definition}

\subsection{Construct}
\subsubsection{Looping}
\begin{definition}[looping]~
    \begin{verbatim}
for item in item1 item2 ...; do 
    cmd
done
\end{verbatim}
\end{definition}

\subsubsection{Conditional}
\begin{definition}[conditional]~
    \begin{verbatim}
if [[ cond1 ]]; then
    cmd1
elif [[ cond2 ]]; then
    cmd2
else
    cmd3
fi

# cond
# ==
# != 
# =~ : unquoted regex  
\end{verbatim}
\end{definition}

\subsubsection{Grouping}
\begin{definition}[grouping]
    \begin{verbatim}
( expr ) # subshell env, var vanish when exit 
\end{verbatim}
\end{definition}

\section{AWK}
\url{https://www.gnu.org/software/gawk/manual/gawk.html}

\begin{definition}[awk]
    The basic function of awk is to search files for lines that contain certain patterns.
\end{definition}

\printbibliography
\end{document}